\section{Chiral Hochschild Cohomology and Bulk--Boundary Correspondence}
\label{sec:chiral_hochschild}
\label{sec:chiral_hochschild_expanded}

Bulk observables admit a geometric interpretation via chiral Hochschild cohomology, developed here through dimensional descent from 2d Poisson sigma models and connection to Khan--Zeng's 3d HT Poisson sigma model \cite{KZ25}.

\subsection{Dimensional Descent: From 2d Poisson Sigma Model to 3d HT}
\subsubsection{Review: 2d Poisson Sigma Model}
Consider a 2d topological field theory on a worldsheet $\Sigma$ with target a Poisson manifold $(M, \pi)$, where $\pi \in \Gamma(\wedge^2 TM)$ is the Poisson bivector. The fields are:
\begin{align}
X &: \Sigma \to M \quad \text{(maps)}, \\
\eta &\in \Omega^1(\Sigma, X^* T^*M) \quad \text{(1-forms on $\Sigma$ valued in the cotangent bundle)}.
\end{align}

The action is:
\begin{equation}
\label{eq:2d_PSM_action}
S_{\text{2d-PSM}}[X, \eta] = \int_\Sigma \left( \langle \eta, dX \rangle + \frac{1}{2} \langle \pi(X), \eta \wedge \eta \rangle \right).
\end{equation}
\subsubsection{Loop Space and Hochschild Cohomology}
Cattaneo--Felder and Kontsevich~\cite{Kon03} observed that the space of fields can be rewritten as:
\begin{equation}
\text{Fields}_{2d} = \text{Maps}(\Sigma, T^*[1]M),
\end{equation}
where $T^*[1]M$ is the $[1]$-shifted cotangent bundle (fiber coordinates $\eta_i$ have degree $+1$).

When $\Sigma = S^1$ (the circle), we have:
\begin{equation}
\text{Fields}|_{S^1} = \text{Maps}(S^1, T^*[1]M) = LT^*[1]M,
\end{equation}
the \textbf{loop space} of $T^*[1]M$.

The fundamental result is:
\begin{theorem}[2d PSM and Hochschild Cohomology;
\ClaimStatusProvedElsewhere]
\label{thm:2d_PSM_Hochschild}
Classical observables in the 2d Poisson sigma model on $S^1$ are identified with Hochschild cochains:
\begin{equation}
\mathcal{O}_{\text{2d-PSM}}(S^1) \cong C^\bullet_{\text{Hoch}}(C^\infty(M), C^\infty(M)),
\end{equation}
where the Hochschild differential $d_{\text{Hoch}}$ corresponds to the BV-BRST differential $Q$ in the sigma model.
\end{theorem}

This is the geometric origin of Hochschild cohomology for Poisson manifolds.
\subsubsection{3d HT Analogue: Chiral Hochschild Cohomology}
Now consider the 3d HT analogue. For a theory on $\R_t \times \C_z$, we replace:
\begin{itemize}
\item Worldsheet $\Sigma$ (2d) $\leadsto$ $\C_z$ (complex curve) + topological direction $\R_t$;
\item Target Poisson manifold $(M, \pi)$ $\leadsto$ Poisson vertex algebra $(V, \{\cdot_\lambda \cdot\})$;
\item Loop space $LM$ $\leadsto$ $\Ran(\C)$, the Ran space of finite subsets of $\C$.
\end{itemize}
\begin{definition}[3d HT Poisson Sigma Model]
\label{def:3d_HT_PSM}
Following Khan--Zeng \cite{KZ25}, the \emph{3d HT Poisson sigma model} associated to a freely generated Poisson vertex algebra $V = \text{Sym}(\C[\phi^i, \partial \phi^i, \ldots])$ has fields:
\begin{align}
\Phi^i &: \R \times \C \to \C \quad \text{(scalar superfields)}, \\
\eta_i &\in \Omega^{0,1}(\C) \otimes \Omega^1(\R) \quad \text{(one-forms in $t$ and $(0,1)$-forms in $z$)}.
\end{align}
The action is:
\begin{equation}
\label{eq:3d_PSM_action}
S_{\text{3d-PSM}} = \int_{\R \times \C} \left( \eta_i (d_t + \dbar) \Phi^i + \frac{1}{2} \Pi^{ij}(\Phi, \partial \Phi, \ldots) \eta_i \wedge \eta_j \right),
\end{equation}
where $\Pi^{ij}(\lambda)$ is the $\lambda$-bracket structure:
\begin{equation}
\{\phi^i_\lambda \phi^j\} = \Pi^{ij}(\lambda) = \sum_{n \geq 0} \frac{\Pi^{ij}_n}{\lambda^{n+1}}.
\end{equation}
\end{definition}

\begin{theorem}[Gauge Invariance $\Leftrightarrow$ $\lambda$-Jacobi;
\ClaimStatusProvedElsewhere]
\label{thm:KZ_gauge_invariance}
(Khan--Zeng \cite{KZ25}, Theorem 2.1) The 3d HT Poisson sigma model \eqref{eq:3d_PSM_action} is gauge-invariant under BRST transformations if and only if the $\lambda$-bracket $\{\cdot_\lambda \cdot\}$ satisfies the $\lambda$-Jacobi identity:
\begin{equation}
\{\phi^i_\lambda \{\phi^j_\mu \phi^k\}\} - \{\phi^j_\mu \{\phi^i_\lambda \phi^k\}\} = \{\{\phi^i_\lambda \phi^j\}_{\lambda + \mu} \phi^k\}.
\end{equation}
\end{theorem}

This establishes a direct link between the PVA axioms (verified in Section \ref{sec:PVA_descent}) and the gauge structure of the 3d HT sigma model.

\subsection{Chiral Hochschild Cochains: Definition and Properties}

\begin{definition}[Chiral Hochschild Cochains]
\label{def:chiral_hochschild}
Let $A$ be a vertex algebra (or PVA). The \emph{chiral Hochschild cochain complex} is:
\begin{equation}
C^\bullet_{\text{ch}}(A, A) = \prod_{n \geq 0} \Hom_{\text{VA}}\big(A^{\otimes n}, A((\lambda_1)) \cdots ((\lambda_{n-1}))\big),
\end{equation}
where $\Hom_{\text{VA}}$ denotes vertex algebra homomorphisms (respecting the translation operator $\partial$ and sesquilinearity).
\end{definition}

\subsubsection{Differential on Chiral Hochschild Cochains}

The differential $d_{\text{ch}}$ on $C^\bullet_{\text{ch}}(A, A)$ has two components:

\begin{enumerate}[label=(\roman*)]
\item \textbf{Internal differential $\delta_Q$}: Induced by the BV-BRST differential $Q$ on $A$:
\begin{equation}
(\delta_Q f)(a_1, \ldots, a_n) = Q(f(a_1, \ldots, a_n)) - \sum_{i=1}^n (-1)^{|f| + |a_1| + \cdots + |a_{i-1}|} f(a_1, \ldots, Qa_i, \ldots, a_n).
\end{equation}

\item \textbf{Hochschild differential $\delta_{\text{Hoch}}$}: Encodes composition via the $\lambda$-bracket:
\begin{multline}
(\delta_{\text{Hoch}} f)(a_0, a_1, \ldots, a_n) = \{{a_0}_\lambda f(a_1, \ldots, a_n)\} \\
+ \sum_{i=1}^{n-1} (-1)^i f(a_1, \ldots, \{{a_i}_{\lambda_i} a_{i+1}\}, \ldots, a_n) + (-1)^n f(a_1, \ldots, a_{n-1}) \{{a_n}_{\lambda_n} \cdot\}.
\end{multline}
\end{enumerate}

The total differential is:
\begin{equation}
d_{\text{ch}} = \delta_Q + \delta_{\text{Hoch}}.
\end{equation}

\begin{lemma}[$d_{\text{ch}}^2 = 0$; \ClaimStatusProvedHere]
\label{lem:dch_squared}
The chiral Hochschild differential satisfies $d_{\text{ch}}^2 = 0$.
\end{lemma}
\begin{proof}
This follows from:
\begin{enumerate}[label=(\roman*)]
\item $\delta_Q^2 = 0$ (since $Q^2 = 0$);
\item $\delta_{\text{Hoch}}^2 = 0$ (by the associativity of the vertex algebra operations (equivalently, the Borcherds identity \cite{Kac98}));
\item $\{\delta_Q, \delta_{\text{Hoch}}\} = 0$ (compatibility: $Q$ is a derivation of the $\lambda$-bracket).
\end{enumerate}
Combining these:
\begin{equation}
d_{\text{ch}}^2 = (\delta_Q + \delta_{\text{Hoch}})^2 = \delta_Q^2 + \{\delta_Q, \delta_{\text{Hoch}}\} + \delta_{\text{Hoch}}^2 = 0.
\end{equation}
\end{proof}

\subsection{Tamarkin's Deformation Theory and $*$-Lie Algebras}

The geometric origin of chiral Hochschild cohomology is Tamarkin's theory of deformations of chiral algebras \cite{Tam02}.

\subsubsection{$*$-Lie Bracket and Chiral Deformations}
\begin{definition}[$*$-Lie Algebra Structure]
\label{def:star_lie}
On the space $C^\bullet_{\text{ch}}(A, A)$, define a $*$-Lie bracket:
\begin{equation}
[f *_\lambda g](a_1, \ldots, a_n) = \{f(a_1, \ldots)_\lambda g(\ldots, a_n)\} - (-1)^{|f||g|} \{g(a_1, \ldots)_\lambda f(\ldots, a_n)\}.
\end{equation}
\end{definition}

This bracket satisfies:
\begin{enumerate}[label=(\roman*)]
\item $[f *_\lambda g] = -(-1)^{(|f|-1)(|g|-1)} [g *_{-\lambda-\partial} f]$ (graded antisymmetry),
\item $[f *_\lambda [g *_\mu h]] - [g *_\mu [f *_\lambda h]] = [[f *_\lambda g] *_{\lambda + \mu} h]$ ($*$-Jacobi identity, which follows from the associativity of the brace operations (Proposition~\ref{prop:geometric-braces-well-defined})).
\end{enumerate}

\begin{theorem}[Tamarkin's Deformation--Quantization;
\ClaimStatusProvedElsewhere]
\label{thm:Tamarkin_deformation}
(Tamarkin \cite{Tam02}) Deformations of a chiral algebra $A$ are controlled by the Maurer--Cartan equation in the $*$-Lie algebra $(C^\bullet_{\text{ch}}(A, A), [*_\lambda \cdot])$:
\begin{equation}
d_{\text{ch}} \alpha + \frac{1}{2} [\alpha *_\lambda \alpha] = 0.
\end{equation}

Solutions $\alpha \in C^1_{\text{ch}}(A, A)$ parametrize infinitesimal deformations of the vertex algebra structure on $A$.
\end{theorem}

\subsection{Bulk--Boundary Identification via Factorization Homology}

The relationship between bulk local operators and chiral Hochschild cochains is as follows.

\subsubsection{Factorization Algebras on the Ran Space}

Recall that bulk local operators in a 3d HT theory form a \emph{factorization algebra} on $\C_z$ (Costello--Gwilliam \cite{CG17}). This means:

\begin{enumerate}[label=(\roman*)]
\item For each open $U \subseteq \C$, we have a cochain complex $\mathcal{O}_{\text{bulk}}(U)$;
\item For disjoint $U_1, \ldots, U_n \subseteq \C$, there are multiplication maps:
\begin{equation}
\mathcal{O}_{\text{bulk}}(U_1) \otimes \cdots \otimes \mathcal{O}_{\text{bulk}}(U_n) \longrightarrow \mathcal{O}_{\text{bulk}}(U_1 \cup \cdots \cup U_n),
\end{equation}
satisfying associativity and locality axioms.
\end{enumerate}

\begin{definition}[Ran Space]
\label{def:Ran_space}
The \emph{Ran space} $\Ran(\C)$ is the colimit
\begin{equation}
\Ran(\C) = \operatorname{colim}_{S \in \mathrm{Fin}^{\mathrm{surj}}} \C^S,
\end{equation}
taken over the category $\mathrm{Fin}^{\mathrm{surj}}$ of non-empty finite sets and surjections, where a surjection $\alpha\colon S\twoheadrightarrow S'$ acts by the diagonal map $\Delta_\alpha\colon \C^{S'}\hookrightarrow \C^{S}$. In particular, $\Ran(\C)$ is the space of finite non-empty subsets of $\C$ (without multiplicity), topologized so that points can collide but not separate.
\end{definition}

The factorization algebra $\mathcal{O}_{\text{bulk}}$ extends to a cosheaf on $\Ran(\C)$: over the open stratum $\Ran_{\leq n}(\C) \setminus \Ran_{\leq n-1}(\C)$ of subsets of cardinality exactly~$n$, it is the $S_n$-equivariant sheaf pulled back from $\mathcal{O}_{\text{bulk}}^{\boxtimes n}$ on the configuration space $\C^n \setminus \Delta$, and the factorization multiplication maps encode the behavior as points collide (moving between strata in the Ran space).

\subsubsection{Factorization Homology and Bulk Observables}

\begin{theorem}[Bulk Observables via Factorization Homology;
\ClaimStatusProvedElsewhere]
\label{thm:bulk_factorization}
(Ayala--Francis \cite{AF15}) Let $\mathcal{A}$ be a factorization algebra on $\C$. Then global observables are computed by factorization homology:
\begin{equation}
\mathcal{O}_{\text{bulk}}(\C) \simeq \int_\C \mathcal{A}.
\end{equation}

For 3d HT theories, this identifies bulk local operators with cochains on the Ran space.
\end{theorem}

\subsubsection{Connection to Chiral Hochschild Cochains}

\begin{theorem}[Bulk $=$ Chiral Hochschild Cochains;
\ClaimStatusProvedHere]
\label{thm:bulk_hochschild}
Assume \textrm{(H1)--(H4)}.
Let $A_{\partial}$ be the boundary chiral algebra of a 3d HT theory (obtained by imposing Dirichlet or Neumann boundary conditions). Then bulk local operators are quasi-isomorphic to chiral Hochschild cochains:
\begin{equation}
\mathcal{O}_{\text{bulk}} \simeq C^\bullet_{\text{ch}}(A_{\partial}, A_{\partial}).
\end{equation}
\end{theorem}

\begin{proof}
We construct an explicit quasi-isomorphism and verify that it intertwines the differentials. The argument has four steps: reduction along $\R$ via local constancy, identification of the resulting complex on $\C$ with chiral Hochschild cochains via the holomorphic weight filtration, matching of differentials, and the quasi-isomorphism verification.

\medskip
\textbf{Step 1: Factorization homology and reduction along $\R$.}

Bulk observables on $\C \times \R$ are computed by factorization homology:
\begin{equation}
\label{eq:bulk-fact-hom}
\mathcal{O}_{\mathrm{bulk}}(\C \times \R) \;\simeq\; \int_{\C \times \R} \mathsf{Obs}.
\end{equation}
Since $\mathsf{Obs}$ is locally constant along $\R$ (Proposition~\ref{prop:prefact-Obs}), factorization homology over the product $\C \times \R$ admits a locally constant factorization along the topological direction \cite{AF15}. Concretely, for any interval $I \subset \R$ and disc $D \subset \C$, the restriction $\mathsf{Obs}(D \times I) \to \mathsf{Obs}(D \times I')$ is a quasi-isomorphism whenever $I \hookrightarrow I'$ is an inclusion of intervals.  The $E_1$--structure along $\R$ is therefore homotopically trivial on global sections, and the factorization homology reduces:
\begin{equation}
\label{eq:R-reduction}
\int_{\C \times \R} \mathsf{Obs} \;\simeq\; \int_{\C} \mathsf{Obs}^{\R},
\end{equation}
where $\mathsf{Obs}^{\R} := \mathsf{Obs}(D \times \R)$ denotes the $\R$--invariant (equivalently, $E_1$--coinvariant) part of the local observable complex, viewed as a holomorphic factorization algebra on $\C$ alone.  At the boundary $\{t = 0\} \times \C$, the restriction map
\begin{equation}
\rho : \mathcal{O}_{\mathrm{bulk}} \longrightarrow A_{\partial}
\end{equation}
identifies $A_{\partial}$ with $\mathsf{Obs}^{\R}|_{\{t=0\}}$, the boundary chiral algebra.

\medskip
\textbf{Step 2: Filtration identification via holomorphic weight.}

The factorization algebra $\mathsf{Obs}^{\R}$ on $\C$ admits a \emph{holomorphic weight filtration} $F^p$ defined by pole order at collision divisors: an element $f \in \int_{\C} \mathsf{Obs}^{\R}$ lies in $F^p$ if, when evaluated on $n$--tuples of operator insertions at points $z_1, \ldots, z_n \in \C$, the resulting meromorphic function has poles of total order at most $p$ along all collision diagonals $\{z_i = z_j\}$.

On the Ran space $\Ran(\C)$, the factorization algebra structure decomposes $\int_{\C} \mathsf{Obs}^{\R}$ into components labeled by the number of insertion points.  An element supported on the stratum $\Ran_n(\C) \setminus \Ran_{n-1}(\C)$ of exactly $n$ distinct points defines a multilinear operation
\begin{equation}
\label{eq:Ran-component}
f_n : A_{\partial}^{\otimes n} \longrightarrow A_{\partial}((\lambda_1)) \cdots ((\lambda_{n-1})),
\end{equation}
parametrized by the spectral parameters $\lambda_i$ (dual to the relative positions $z_i - z_n$).  The collection $\{f_n\}_{n \geq 0}$ is precisely an element of the chiral Hochschild cochain complex $C^\bullet_{\mathrm{ch}}(A_{\partial}, A_{\partial})$ (Definition~\ref{def:chiral_hochschild}).

The holomorphic weight filtration on the bulk side matches the \emph{Hochschild filtration} on $C^\bullet_{\mathrm{ch}}(A_{\partial}, A_{\partial})$: filtration degree $p$ on the bulk corresponds to cochains supported on $\Ran_{\leq p}(\C)$, i.e., cochains of at most $p$ inputs.  This is because each additional insertion point introduces a new collision divisor contributing to the pole order.

\medskip
\textbf{Step 3: Differential identification.}

The total differential $d_{\mathrm{bulk}}$ on $\int_{\C} \mathsf{Obs}^{\R}$ has two components, which we identify with the two components of the chiral Hochschild differential $d_{\mathrm{ch}} = \delta_Q + \delta_{\mathrm{Hoch}}$.

\begin{enumerate}[label=(\roman*)]
\item \emph{Internal differential $\delta_Q$:} The BV-BRST differential $Q = m_1$ acts on each factor of $A_{\partial}$ separately.  On a cochain $f_n$, this gives
\[
(\delta_Q f_n)(a_1, \ldots, a_n) = Q(f_n(a_1, \ldots, a_n)) - \sum_{i=1}^n (-1)^{|f| + |a_1| + \cdots + |a_{i-1}|} f_n(a_1, \ldots, Qa_i, \ldots, a_n).
\]
This is the graded commutator of $Q$ with $f_n$, corresponding to the linearized BRST operator on bulk observables.

\item \emph{Hochschild differential $\delta_{\mathrm{Hoch}}$:} The factorization product (gluing bulk operators as insertion points collide) is controlled by the codimension-1 boundary faces of $\FM_k(\C)$.  For an $n$--cochain $f_n$, the Hochschild differential inserts $m_2$ at adjacent positions:
\[
(\delta_{\mathrm{Hoch}} f_n)(a_1, \ldots, a_{n+1}) = \sum_{i=1}^{n} \pm\, f_n(a_1, \ldots, m_2(a_i, a_{i+1})\big|_{\Lambda_{\{i,i+1\}}}, \ldots, a_{n+1}),
\]
together with the two end terms where $a_1$ or $a_{n+1}$ acts on the output via the $\lambda$--bracket.  These are precisely the boundary face maps: the divisor $D_{\{i,i+1\}} \subset \partial \FM_{n+1}(\C)$ where points $z_i$ and $z_{i+1}$ collide corresponds to composing $f_n$ with $m_2$ at the $i$th slot, with spectral parameter specialization $\Lambda_{\{i,i+1\}} = \lambda_i + \lambda_{i+1}$ implementing the block fusion rule (spectral substitution).
\end{enumerate}

The two components together give
\begin{equation}
d_{\mathrm{bulk}} = \delta_Q + \delta_{\mathrm{Hoch}} = d_{\mathrm{ch}},
\end{equation}
so the map $\int_{\C} \mathsf{Obs}^{\R} \to C^\bullet_{\mathrm{ch}}(A_{\partial}, A_{\partial})$ defined by the Ran space decomposition~\eqref{eq:Ran-component} is a map of cochain complexes.

\medskip
\textbf{Step 4: Quasi-isomorphism.}

We verify that the map
\begin{equation}
\Phi : \int_{\C} \mathsf{Obs}^{\R} \longrightarrow C^\bullet_{\mathrm{ch}}(A_{\partial}, A_{\partial})
\end{equation}
is a quasi-isomorphism by a filtration argument.  Both sides carry the holomorphic weight/Hochschild filtration, and $\Phi$ is a filtered map.  On the associated graded, the factorization algebra $\gr^F\!\int_{\C} \mathsf{Obs}^{\R}$ splits: at each filtration level $p$, the stratum $\Ran_p(\C) \setminus \Ran_{p-1}(\C)$ contributes $\Hom(A_{\partial}^{\otimes p}, A_{\partial}((\lambda_1))\cdots((\lambda_{p-1})))$ with differential given purely by $\delta_Q$ (since $\delta_{\mathrm{Hoch}}$ raises filtration degree).  The associated graded of $C^\bullet_{\mathrm{ch}}(A_{\partial}, A_{\partial})$ has the same description.  Hence $\gr(\Phi)$ is an isomorphism, and by completeness of the filtration and the standard spectral sequence comparison theorem, $\Phi$ is a quasi-isomorphism:
\begin{equation}
C^\bullet_{\mathrm{ch}}(A_{\partial}, A_{\partial}) \;\simeq\; \int_{\C} \mathsf{Obs}^{\R} \;\simeq\; \mathcal{O}_{\mathrm{bulk}}(\C \times \R),
\end{equation}
as dg-objects, where $C^\bullet_{\mathrm{ch}}$ is the chiral Hochschild cochain complex.
\end{proof}

\subsection{Comparison with CDG Boundary Chiral Algebras}

Our construction recovers and extends the boundary chiral algebra framework of Costello--Dimofte--Gaiotto \cite{CDG20}:

\begin{theorem}[CDG Compatibility; \ClaimStatusProvedHere]
\label{thm:CDG_compatibility}
Assume \textrm{(H1)--(H4)}.
For 3d $\mathcal{N}=2$ theories with holomorphic-topological twist:
\begin{enumerate}[label=(\roman*)]
\item \textbf{Bulk is commutative PVA}: $H^\bullet(\mathcal{O}_{\text{bulk}}, Q)$ is a commutative Poisson vertex algebra with $(-1)$-shifted $\lambda$-bracket (Theorem \ref{thm:cohomology_PVA});

\item \textbf{Boundary is vertex algebra}: $A_{\partial}$ is a vertex algebra (may or may not have Virasoro);

\item \textbf{Bulk-to-boundary map}: There is a canonical morphism of PVAs
\begin{equation}
\beta : H^\bullet(\mathcal{O}_{\text{bulk}}, Q) \longrightarrow Z(A_{\partial}),
\end{equation}
where $Z(A_{\partial})$ is the center of $A_{\partial}$ (elements commuting with all $\lambda$-brackets);

\item \textbf{Boundary is bulk module}: $A_{\partial}$ is a module for the bulk PVA $H^\bullet(\mathcal{O}_{\text{bulk}}, Q)$ via the map $\beta$.
\end{enumerate}
\end{theorem}

\begin{proof}
Parts (i) and (ii) are established in \cite{CDG20}, Section 4. Part (iii) follows from our bulk--Hochschild identification (Theorem \ref{thm:bulk_hochschild}): the center $Z(A_{\partial})$ corresponds to $HH^0(A_{\partial}, A_{\partial})$, which is the degree-0 component of chiral Hochschild cohomology. Part (iv) is the module structure induced by the bulk-to-boundary restriction map.
\end{proof}

\subsection{Chiral Hochschild cochains as a brace algebra on $\FM(\C)\times\Conf(\R)$}
\label{sec:chiral-Hochschild}

\subsubsection{Operadic cochains for colored operads and the brace structure}
\label{subsec:operadic-cochains}
Let $O$ be a colored dg--operad and $A$ an $O$--algebra in a symmetric monoidal dg--category $(\mathcal C,\otimes)$.
Write $O^\vee$ for the bar cooperad of $O$ (cofree conilpotent cooperad on $s^{-1}\overline O$).
We adopt the standard operadic cochain complex
\[
\mathrm{C}^\bullet_O(A) \;:=\; \underline{\mathrm{Hom}}_{\text{cooperad}} \bigl( O^\vee,\, \mathrm{End}_A \bigr),
\]
equipped with the differential induced by the coderivation on $O^\vee$ and the internal differential of $\mathrm{End}_A$.
When $O = \OHT$ is our Swiss--cheese chain operad, we denote
\[
\mathrm{C}^\bullet_{\mathrm{ch,top}}(A)\;:=\;\mathrm{C}^\bullet_{\OHT}(A)
\quad\text{(the \emph{chiral Hochschild cochains})}.
\]

\begin{definition}[Brace operations]
\label{def:brace-ops}
The bar cooperad $O^\vee$ carries a canonical \emph{brace} structure; by precomposition,
$\mathrm{C}^\bullet_O(A)$ acquires multilinear operations
\[
\{-\}\{-,\ldots,-\}\colon \mathrm{C}^\bullet_O(A)\otimes \mathrm{C}^\bullet_O(A)^{\otimes m}\longrightarrow \mathrm{C}^\bullet_O(A),
\]
given by grafting trees in $O^\vee$ (see \cite{GJ94,Vor99}).
The resulting algebraic structure on $\mathrm{C}^\bullet_O(A)$ is a homotopy Gerstenhaber (brace) algebra; its induced \emph{Gerstenhaber bracket} is
\[
[f,g] := f\{g\} - (-1)^{(\deg f-1)(\deg g-1)}\, g\{f\}.
\]
\end{definition}

\begin{remark}[Degree conventions]
We use cohomological grading. For $O=\OHT$ the internal degrees reflect the form--degrees on $\FM(\C)$ and the $E_1$--degrees on $\Conf(\R)$; the brace degree shifts match the usual Gerstenhaber sign rule.
\end{remark}

\subsubsection{Geometric brace model on $\FM(\C)\times\Conf(\R)$}
\label{subsec:geometric-brace}
A concrete model of the brace operations uses configuration spaces.
Let $\Conf_k(\R)$ denote ordered configurations $(t_1<\cdots<t_k)$ of $k$ points in~$\R$ and set
\[
\mathbb{X}_{k,m} \;:=\; \FM_k(\C)\times \Conf_m(\R).
\]
Operations in $\OHT$ at arity $(k,m)$ are represented by chains on $\mathbb{X}_{k,m}$ decorated by edge--length parameters.
Given $f\in \mathrm{C}^\bullet_{\mathrm{ch,top}}(A)$ and $g_1,\dots,g_m\in \mathrm{C}^\bullet_{\mathrm{ch,top}}(A)$, we define
\[
f\{g_1,\dots,g_m\} := \int_{\mathbb{X}} \;\;\Bigl(\text{pullback of kernels from } \mathbb{X}_{k,m'} \Bigr)\;\circ\;\Bigl(\text{insertions via trees along } \FM\times \Conf\Bigr),
\]
where the integral is over the appropriate fiber product of configuration spaces determined by the brace tree (see Appendix~\ref{app:brace-signs} for precise sign and degree formulas).

\begin{proposition}[Well-definedness and functoriality;
\ClaimStatusProvedHere]
\label{prop:geometric-braces-well-defined}
Under the analytic hypotheses (H1)--(H4), the geometric brace operations on $\mathrm{C}^\bullet_{\mathrm{ch,top}}(A)$ are well-defined, strictly multilinear, and satisfy the brace identities. The Gerstenhaber bracket obtained from these braces coincides with the one induced by the cooperad $O^\vee$.
\end{proposition}

\begin{proof}
Integrability follows from Hypothesis~(H2): the weight forms have at most logarithmic singularities along the boundary divisors of the FM compactification, and such forms are integrable on compact manifolds with corners. The brace identities follow from Stokes' theorem on $\FM_k(\C) \times \Conf_m(\R)$: the boundary strata reproduce the algebraic brace relations, with the Arnold cancellations at codimension-2 corners ensuring that no spurious terms arise (see Appendix~\ref{app:brace-signs} for the detailed sign verification). Compatibility with the cooperad differential is standard.
\end{proof}

\subsubsection{The \texorpdfstring{$\ast$}{*}--Lie structure and deformation theory}
\label{subsec:star-Lie}
Let $(\mathrm{C}^\bullet_{\mathrm{ch,top}}(A),\{-\}\{-\})$ be the brace algebra constructed above.

\begin{definition}[$\ast$--Lie algebra]
\label{def:star-Lie}
Define the shifted Lie bracket on $\mathrm{C}^\bullet_{\mathrm{ch,top}}(A)[1]$ by the Gerstenhaber bracket $[-,-]$ associated to braces. We call $(\mathrm{C}^\bullet_{\mathrm{ch,top}}(A)[1],[-,-])$ the \emph{chiral $\ast$--Lie algebra} of $A$.
\end{definition}

\begin{theorem}[Deformations and Maurer--Cartan; \ClaimStatusProvedHere]
\label{thm:MC-deformations}
Let $A$ be a $\OHT$--algebra, and consider formal deformations over $k[\![\hbar]\!]$. Then there is a bijection between:
\begin{enumerate}[label=(\alph*)]
  \item deformation classes of $A$ as a $W(\SCchtop)$--algebra over $k[\![\hbar]\!]$ up to gauge; and
  \item gauge--equivalence classes of solutions to the Maurer--Cartan equation
    \(
      d\alpha + \tfrac{1}{2}[\alpha,\alpha]=0
    \)
    in $\bigl(\mathrm{C}^\bullet_{\mathrm{ch,top}}(A)[1]\otimes \hbar k[\![\hbar]\!],[-,-]\bigr)$.
\end{enumerate}
\end{theorem}

\begin{proof}
This is the standard operadic deformation theory: deformations are classified by twisting morphisms from the bar cooperad into $\mathrm{End}_A$, and a twisting morphism is equivalent to an MC element in the convolution Lie algebra $\mathrm{Hom}(O^\vee,\mathrm{End}_A)$; see \cite{GK94,Val07}. The brace model identifies this Lie algebra with $\mathrm{C}^\bullet_{\mathrm{ch,top}}(A)[1]$.
\end{proof}

\subsubsection{Bulk observables $\simeq$ chiral Hochschild cochains}
\label{subsec:bulk-equals-CHC}
Let $\mathsf{Obs}$ be the BV observable prefactorization algebra of the HT theory, and $A=\mathsf{Obs}(D\times I)$ the local complex on a small rectangle.

\begin{theorem}[Bulk--cochains identification; \ClaimStatusProvedHere]
\label{thm:bulk-CHC}
Assume \textrm{(H1)--(H4)}.
There is a canonical filtered quasi--isomorphism
\[
\mathrm{C}^\bullet_{\mathrm{ch,top}}(A)\;\simeq\;\mathsf{Obs}_{\mathrm{bulk}}(\mathbb{B}^3)
\]
between chiral Hochschild cochains of the local algebra $A$ and the complex of bulk local observables supported in a small $3$--ball, with filtration by holomorphic weight. On associated graded, this identification reduces to the BD--chiral Hochschild complex on the closed color tensored with the $E_1$ Hochschild complex on the open color.
\end{theorem}

\begin{proof}
The argument proceeds in three steps: an operad-of-shapes recognition that provides a common model for both sides, a filtration analysis via holomorphic weight, and a well-definedness verification using the analytic hypotheses.

\medskip
\textbf{Step 1: Operad-of-shapes recognition.}

Decompose the small $3$--ball $\mathbb{B}^3 \subset \C \times \R$ into rectangles $D_i \times I_j$ (discs in $\C$ times intervals in $\R$).  On each rectangle, $\mathsf{Obs}(D_i \times I_j) \simeq A$ by local constancy (Proposition~\ref{prop:prefact-Obs}).  The factorization homology $\int_{\mathbb{B}^3} \mathsf{Obs}$ is computed as a homotopy colimit over all such decompositions; the indexing category is the \emph{operad of shapes} $\mathbb{X}_{k,m} = \FM_k(\C) \times \Conf_m(\R)$, where $k$ holomorphic insertion points in $\C$ and $m$ ordered topological insertion points in $\R$ parametrize the ways that rectangles embed in $\mathbb{B}^3$.

Concretely, an element of $\mathsf{Obs}_{\mathrm{bulk}}(\mathbb{B}^3)$ is specified by a compatible family of multilinear operations
\[
A^{\otimes k} \otimes A^{\otimes m} \longrightarrow A,
\]
parametrized by points in $\mathbb{X}_{k,m}$, i.e., by integration of propagator kernels over $\FM_k(\C) \times \Conf_m(\R)$.  By the Swiss--cheese recognition (Theorem~\ref{conj:recognition}), these operations assemble into a $W(\SCchtop)$--algebra, and the totality of such operations is precisely the operadic cochain complex
\[
\mathrm{C}^\bullet_{\mathrm{ch,top}}(A) = \underline{\mathrm{Hom}}_{\text{cooperad}}\bigl((\OHT)^\vee, \mathrm{End}_A\bigr).
\]
This gives the map $\Phi \colon \mathrm{C}^\bullet_{\mathrm{ch,top}}(A) \to \mathsf{Obs}_{\mathrm{bulk}}(\mathbb{B}^3)$: a cochain $f$ of arity $(k,m)$ determines an observable by integrating $f$ against the propagator kernels over $\mathbb{X}_{k,m}$.  In the reverse direction, restricting a bulk observable to the configuration-space strata recovers the cochain components.

\medskip
\textbf{Step 2: Holomorphic weight filtration.}

Equip both sides with the \emph{holomorphic weight filtration} $F^p$, graded by total pole order at collision divisors in $\FM_k(\C)$.

On the cochain side, define $F^p \mathrm{C}^\bullet_{\mathrm{ch,top}}(A) \subset \mathrm{C}^\bullet_{\mathrm{ch,top}}(A)$ as the subcomplex of cochains whose kernels, when restricted to $\FM_k(\C)$, have poles of total order at most $p$ along the union of collision divisors $\bigcup_{|S| \geq 2} D_S$.

On the observable side, define $F^p \mathsf{Obs}_{\mathrm{bulk}}(\mathbb{B}^3)$ by the analogous pole-order condition on the Feynman-diagrammatic integrands.

The map $\Phi$ is filtered: it preserves pole order since it is defined by integration over the \emph{same} configuration-space kernels.  On the associated graded, the factorization structure splits by the product decomposition $\mathbb{X}_{k,m} = \FM_k(\C) \times \Conf_m(\R)$.  At each graded piece:
\begin{enumerate}[label=(\roman*)]
\item The \emph{chiral component} (varying $k$, from $\FM_k(\C)$) contributes the BD--chiral Hochschild complex $\mathrm{C}^\bullet_{\mathrm{BD\text{-}ch}}(A_{\mathsf{ch}})$, where $A_{\mathsf{ch}}$ is the closed-color (holomorphic) part of the $W(\SCchtop)$--algebra.  The differential on this piece is the chiral Hochschild differential with $\delta_Q$ from the internal BV-BRST operator and $\delta_{\mathrm{Hoch}}$ from the boundary faces of $\FM_k(\C)$ encoding OPE collisions.
\item The \emph{topological component} (varying $m$, from $\Conf_m(\R)$) contributes the $E_1$ Hochschild complex $\mathrm{C}^\bullet_{E_1}(A_{\mathsf{top}})$, where $A_{\mathsf{top}}$ is the open-color (topological) part.  The differential encodes the boundary faces of $\Conf_m(\R)$ (merging adjacent intervals).
\end{enumerate}
The splitting on associated graded is
\[
\gr^F \mathrm{C}^\bullet_{\mathrm{ch,top}}(A) \;\cong\; \mathrm{C}^\bullet_{\mathrm{BD\text{-}ch}}(A_{\mathsf{ch}}) \;\widehat\otimes\; \mathrm{C}^\bullet_{E_1}(A_{\mathsf{top}}),
\]
and the same decomposition holds for $\gr^F \mathsf{Obs}_{\mathrm{bulk}}(\mathbb{B}^3)$ by the factorization product structure on disjoint rectangles.  Hence $\gr(\Phi)$ is an isomorphism.

\medskip
\textbf{Step 3: Well-definedness and integrability.}

The integrals defining $\Phi$ converge by Hypothesis~(H2) (Propagator Regularity): the propagator kernels are meromorphic in $\C$ with at most simple poles at collision points, and exponentially decaying along $\R$.  After pullback to the FM compactification $\FM_k(\C) \times \Conf_m(\R)$, the resulting forms have at most logarithmic singularities along the boundary divisors $D_S$.  Since $\FM_k(\C)$ is a compact manifold with corners, forms with logarithmic singularities are integrable --- the key estimate is
\[
\int_0^\epsilon |\log \varepsilon_S|^N \, \varepsilon_S^{a-1} \, d\varepsilon_S < \infty \quad \text{for } a > 0, \; N \geq 0,
\]
where $\varepsilon_S$ is the radial parameter measuring the size of the cluster $S$ (Definition~\ref{def:FM_coords}).  The exponent $a > 0$ is guaranteed by the simple-pole condition in (H2); higher-loop contributions are controlled by (H1) (one-loop finiteness ensures no accumulation of poles).

The filtration is complete and Hausdorff on both sides (pole order is bounded below and the complexes are exhausted by finite filtration levels).  By the spectral sequence comparison theorem, the isomorphism on $\gr$ lifts to a quasi-isomorphism on the filtered complexes:
\[
\Phi \colon \mathrm{C}^\bullet_{\mathrm{ch,top}}(A) \;\xrightarrow{\;\simeq\;}\; \mathsf{Obs}_{\mathrm{bulk}}(\mathbb{B}^3). \qedhere
\]
\end{proof}

\begin{corollary}[Obstruction theory; \ClaimStatusProvedHere]
\label{cor:obstructions}
Obstructions to first--order deformations lie in $H^2(\mathrm{C}^\bullet_{\mathrm{ch,top}}(A))$ and equivalence classes of infinitesimal deformations are parameterized by $H^1$. Higher obstructions live in the Massey--product tower controlled by the brace operations.
\end{corollary}

\subsubsection{Filtration and PVA on associated graded}
\label{subsec:filtration-PVA}
Let $F^\bullet$ be the holomorphic weight filtration on $\mathrm{C}^\bullet_{\mathrm{ch,top}}(A)$. Then
\[
\mathrm{gr}^F\, \mathrm{C}^\bullet_{\mathrm{ch,top}}(A)
\;\cong\;
\mathrm{C}^\bullet_{\mathrm{BD\text{-}ch}}(A_{\mathsf{ch}})\;\widehat\otimes\; \mathrm{C}^\bullet_{E_1}(A_{\mathsf{top}}),
\]
and the induced bracket on $\mathrm{H}^\bullet$ coincides with the $(-1)$--shifted Poisson vertex bracket of Section~\ref{sec:Ainfty-to-PVA}.

\begin{theorem}[Hochschild bridge to Volume~I at genus~$0$;
\ClaimStatusProvedHere]
\label{thm:hochschild-bridge-genus0}
Assume \textrm{(H1)--(H4)}.  Let $\cA$ be the BV-BRST complex of a
3d HT theory on $\C \times \R$, and assume $\cA$ is chiral Koszul
(Volume~I, Theorem~B).  Then there is a canonical filtered
quasi-isomorphism
\[
\mathrm{C}^\bullet_{\mathrm{ch,top}}(\cA)
\;\simeq\;
\ChirHoch^\bullet(\cA),
\]
where $\ChirHoch^\bullet(\cA)$ is the chiral Hochschild complex of
Volume~I on the curve $X = \C$.  On the Koszul locus, this complex
is polynomial in $r$~generators of degrees $m_i + 1$
(Volume~I, Theorem~H).
\end{theorem}

\begin{proof}
Three steps.

\textbf{Step~1} (Translation invariance and local constancy).
The 3d HT theory on $\C \times \R$ is translation-equivariant in
the holomorphic direction: the action, propagator, and BV-BRST
differential are all invariant under $z \mapsto z + c$.
Consequently, the bulk observables
$\mathrm{Obs}_{\mathrm{bulk}}$ form a \emph{locally constant}
factorization algebra on~$\C$ (after integrating out the
$\R$-direction).  The chiral Hochschild cochains
$\ChirHoch^\bullet(\cA)$, defined by Volume~I via configuration
spaces $C_{n+2}(\C)$ with logarithmic forms, inherit the same
translation equivariance and therefore also form a locally
constant factorization algebra on~$\C$.

\textbf{Step~2} (Contractibility and localization).
Since $\C$ is contractible, every locally constant factorization
algebra on $\C$ is determined up to quasi-isomorphism by its
stalk at any point $z_0 \in \C$
\cite[Theorem~5.5.1.1]{CG17}.  The stalk of
$\ChirHoch^\bullet(\cA)$ at $z_0$ is computed by restricting to a
small disc $D \ni z_0$:
\[
\ChirHoch^\bullet(\cA)|_{z_0}
\;\simeq\;
\bigoplus_{n \geq 0}
  \Gamma\!\bigl(
    C_{n+2}(D),\,
    j_\ast j^\ast \cA^{\otimes(n+2)}
    \otimes \Omega^n_{\log}
  \bigr).
\]
This is exactly the local complex
$\mathrm{C}^\bullet_{\mathrm{ch,top}}(\cA)$ of
Theorem~\ref{thm:bulk-CHC}, restricted from $\FM_{n+2}(\C)$ to
$\FM_{n+2}(D)$.  Since $D \hookrightarrow \C$ induces a homotopy
equivalence on ordered configuration spaces (both being
configuration spaces in a contractible domain with the same
fundamental group of the complement), the restriction is a
quasi-isomorphism.

\textbf{Step~3} (Filtration and polynomial structure).
Both sides carry the holomorphic weight filtration $F^p$ counting
pole order along collision divisors in $\FM_k(\C)$.  The
localization map of Step~2 preserves pole order (restriction does
not change singularity type), so it is a \emph{filtered}
quasi-isomorphism.  On the Koszul locus, Volume~I's Theorem~H
gives $\ChirHoch^\bullet(\cA) \cong \C[\Theta_1, \ldots,
\Theta_r]$ with $\deg \Theta_i = m_i + 1$; this polynomial
structure transports across the identification.
\end{proof}

\begin{remark}[Connes periodicity]
Under the identification of
Theorem~\ref{thm:hochschild-bridge-genus0}, the Connes periodicity
operator $S$ on $\ChirHoch^\bullet$ (shifting cyclic degree by
$-2$) corresponds to the holomorphic weight shift
$F^p \to F^{p-1}$ on
$\mathrm{C}^\bullet_{\mathrm{ch,top}}(\cA)$.  This follows from
the fact that $S$ acts on configuration spaces by ``forgetting one
insertion point,'' which lowers the number of collision divisors
and hence the pole order by one unit in each factor.  We record
this as a consequence rather than an ingredient of the proof.
\end{remark}

\begin{theorem}[Hochschild bridge at genus~$g \geq 1$;
\ClaimStatusProvedHere]
\label{thm:hochschild-bridge-higher-genus}
Assume \textrm{(H1)--(H4)}.  Let $\cA$ be a chiral Koszul algebra
on a smooth projective curve $\Sigma_g$ of genus~$g \geq 1$.
Then there is a canonical filtered quasi-isomorphism of complexes
over $\overline{\mathcal{M}}_g$:
\begin{equation}\label{eq:hochschild-bridge-genus-g}
\mathrm{C}^\bullet_{\mathrm{ch,top}}(\cA;\, \Sigma_g)
\;\simeq\;
\ChirHoch^\bullet(\cA;\, \Sigma_g),
\end{equation}
where the left side is the genus-$g$ chiral Hochschild cochains
(the curved complex with $d^2 = \kappa(\cA) \cdot \omega_g$, as
in Proposition~\textup{\ref{prop:curved-R-factorization}}), and
the right side is Volume~I's chiral Hochschild complex on
$\Sigma_g$.

The quasi-isomorphism has the following structure:
\begin{enumerate}[label=\textup{(\roman*)}]
\item \emph{Stalk identification.}
  On each fiber over a point $[\Sigma_g] \in \mathcal{M}_g$,
  the stalk of $\ChirHoch^\bullet(\cA;\, \Sigma_g)$ at any
  $z_0 \in \Sigma_g$ is quasi-isomorphic to
  $\mathrm{C}^\bullet_{\mathrm{ch,top}}(\cA)$ (the local
  complex of Theorem~\textup{\ref{thm:hochschild-bridge-genus0}}).

\item \emph{Monodromy correction.}
  The global sections on $\Sigma_g$ differ from the local stalk
  by the period-matrix corrections of Volume~I's genus tower:
  \[
  \Gamma(\Sigma_g,\, \ChirHoch^\bullet(\cA))
  \;\simeq\;
  \mathrm{C}^\bullet_{\mathrm{ch,top}}(\cA)
  \;\widehat\otimes\;
  \Lambda^\bullet\bigl(H^1(\Sigma_g)\bigr),
  \]
  where the exterior algebra factor records the
  $2g$~independent monodromies of the factorization algebra
  around the cycles of~$\Sigma_g$.

\item \emph{Curvature.}
  The curved differential on the left satisfies
  $d^2 = \kappa(\cA) \cdot \omega_g$
  (Proposition~\textup{\ref{prop:curved-R-factorization}}).
  Under the identification, this curvature maps to the
  genus-$g$ obstruction class $[\Theta_g] \in
  H^2_{\mathrm{cyc}}(\ChirHoch^\bullet(\cA))$ of Volume~I
  (Theorem~D\textsubscript{scal}): the physical curvature IS the
  algebraic obstruction to extending the Hochschild complex
  from genus~$0$ to genus~$g$.

\item \emph{Complementarity.}
  For a Koszul pair $(\cA, \cA^!)$, the curvatures cancel:
  $\kappa(\cA) + \kappa(\cA^!) = 0$
  (Volume~I, Theorem~C).  Hence the tensor product
  $\ChirHoch^\bullet(\cA) \otimes \ChirHoch^\bullet(\cA^!)$
  is uncurved at every genus, and the combined Hochschild
  complex extends to a flat family over
  $\overline{\mathcal{M}}_g$.
\end{enumerate}
\end{theorem}

\begin{proof}
\textbf{Step~1} (Stalk identification).
Fix $[\Sigma_g] \in \mathcal{M}_g$ and $z_0 \in \Sigma_g$.
In a small disc $D \ni z_0$, the curve is biholomorphic to
$\C$ (uniformization), so the local computation on $D \times I$
is identical to the genus-$0$ case.  By
Theorem~\ref{thm:hochschild-bridge-genus0} applied to the
restriction $\cA|_D$, the stalk of $\ChirHoch^\bullet(\cA)$ at
$z_0$ is quasi-isomorphic to the local complex
$\mathrm{C}^\bullet_{\mathrm{ch,top}}(\cA)$.

\textbf{Step~2} (Global assembly via factorization descent).
Cover $\Sigma_g$ by contractible opens $\{U_\alpha\}$.  On each
$U_\alpha$, the factorization algebra
$\ChirHoch^\bullet(\cA)|_{U_\alpha}$ is locally constant
(by translation invariance in local coordinates) and hence
determined by its stalk (Step~1).  The global sections are
computed by the \v{C}ech--Thom--Sullivan totalization of the
factorization cosheaf
\cite[Chapter~6]{CG17}:
\begin{equation}\label{eq:cech-ts}
\Gamma(\Sigma_g,\, \ChirHoch^\bullet(\cA))
\;\simeq\;
\mathrm{Tot}\!\bigl(
  \check{C}^\bullet(\{U_\alpha\},\,
  \ChirHoch^\bullet(\cA))
\bigr).
\end{equation}
For a locally constant factorization algebra on $\Sigma_g$,
this totalization is computed by the representation of
$\pi_1(\Sigma_g)$ on the stalk.  The result is the stalk
tensored with $\Lambda^\bullet(H^1(\Sigma_g))$: each
$1$-cycle contributes an exterior generator recording the
monodromy of the factorization algebra around that cycle.
This is the standard computation for locally constant sheaves
on surfaces (equivalently: $\Gamma(\Sigma_g,\,
\underline{V}) \simeq V \otimes
H^\bullet(\Sigma_g)$ for a locally constant sheaf with
stalk~$V$ and trivial monodromy, and the exterior algebra
$\Lambda^\bullet(H^1(\Sigma_g))$ records the monodromy
data via the exponential map
$H^1(\Sigma_g;\, \mathrm{End}(V)) \to
\mathrm{Aut}(V)^{2g}$).

\textbf{Step~3} (Curvature identification).
The curved bar differential on $\barB^{(g)}(\cA)$ satisfies
$d_{\mathrm{fib}}^2 = \kappa(\cA) \cdot \omega_g$
(Proposition~\ref{prop:curved-R-factorization}), where
$\kappa(\cA)$ is the modular characteristic (Volume~I,
Theorem~D) and $\omega_g$ is the Arakelov $(1,1)$-form.
Volume~I identifies $\kappa(\cA) \cdot \omega_g$ with the
genus-$g$ obstruction class in cyclic cohomology
$H^2_{\mathrm{cyc}}(\cA, \cA)$ via the Chern--Weil map
from the Hodge bundle $\mathbb{E} \to \overline{\mathcal{M}}_g$
(Theorem~D\textsubscript{scal}).
Under the stalk identification of Step~1, the physical curvature
(failure of $d^2 = 0$ on the fiberwise bar complex) maps to the
algebraic obstruction (failure of the Hochschild complex to
extend as a dg complex from genus~$0$ to genus~$g$).

\textbf{Step~4} (Complementarity).
For a Koszul pair $(\cA, \cA^!)$, Volume~I's Theorem~C gives
$\kappa(\cA) + \kappa(\cA^!) = 0$.  The tensor product
$\ChirHoch^\bullet(\cA) \otimes \ChirHoch^\bullet(\cA^!)$
therefore has curvature
$(\kappa(\cA) + \kappa(\cA^!)) \cdot \omega_g = 0$, and the
combined complex is a genuine dg complex (not curved) at every
genus.  The Hochschild bridge for the combined system is
therefore a quasi-isomorphism of honest (uncurved) dg algebras
over $\overline{\mathcal{M}}_g$.
\end{proof}

\begin{remark}[The full Hochschild bridge]
Theorems~\ref{thm:hochschild-bridge-genus0}
and~\ref{thm:hochschild-bridge-higher-genus} together constitute
the \emph{complete Hochschild bridge} between the two volumes:
the physical bulk observables of the 3d HT theory, computed via
BV-BRST (Volume~II), ARE the chiral Hochschild cochains of
Volume~I's categorical logarithm, at all genera.  The genus-$0$
case is a stalk computation.  The genus-$g$ extension is a
factorization descent that introduces monodromy
($\Lambda^\bullet(H^1(\Sigma_g))$) and curvature
($\kappa(\cA) \cdot \omega_g$) --- both already computed in
Volume~I.

With this bridge closed, all five of Volume~I's main theorems
(A through H) have their physical counterparts established in
Volume~II:
\begin{center}
\renewcommand{\arraystretch}{1.2}
\begin{tabular}{lll}
\textbf{Theorem} & \textbf{Vol~I (algebra)} &
  \textbf{Vol~II (physics)} \\
\hline
A (bar-cobar) & Factorization coalgebra &
  Swiss-cheese algebra on $\C \times \R$ \\
B (inversion) & Quillen equivalence &
  Line operators $\simeq$ $\cA^!$-modules \\
C (complementarity) & $\kappa + \kappa^! = 0$ &
  Anomaly cancellation for Koszul pairs \\
D (leading coeff.) & $\kappa(\cA) \cdot \omega_g$ &
  Curved HT at genus $g$ \\
H (Hochschild) & Polynomial ring &
  Bulk $\simeq$ chiral Hochschild
\end{tabular}
\end{center}
\end{remark}

\begin{corollary}[Bridge to Volume~I Theorem~H; \ClaimStatusConjectured]
\label{cor:hochschild-bridge}
When specialized to a chiral algebra $\cA$ on a curve $X$ with
trivial topological direction ($\R = \mathrm{pt}$), the chiral
Hochschild complex of this volume recovers the chiral Hochschild
complex $\ChirHoch^*(\cA)$ of Volume~I\@.  In particular, for a
principal $\mathcal{W}$-algebra $\mathcal{W}^k(\fg)$ at generic
level, the result is polynomial:
$\ChirHoch^*(\mathcal{W}^k(\fg)) \cong
\mathbb{C}[\Theta_1, \ldots, \Theta_r]$.
\end{corollary}

\begin{proof}[Proof sketch]
When $\R = \mathrm{pt}$, the $\SCchtop$ structure degenerates: the
open color (topological direction) contributes only a one-dimensional
$E_1$-algebra, so the Swiss-cheese cochain complex
$\mathrm{C}^\bullet_{\mathrm{ch,top}}(\cA)$ reduces to the purely
holomorphic factorization cochain complex on~$X$ alone.  Concretely,
the configuration space factor $\Conf_m(\R)$ collapses to a point for
each $m$, and the brace operations (Definition~\ref{def:brace-ops})
reduce to classical Gerstenhaber braces on the closed-color
(holomorphic) sector.  The resulting complex is the BD-chiral
Hochschild complex $\mathrm{C}^\bullet_{\mathrm{BD\text{-}ch}}(\cA)$,
which coincides with Volume~I's $\ChirHoch^\bullet(\cA)$ by
Theorem~\ref{thm:hochschild-bridge-genus0}.

For $\cA = \mathcal{W}^k(\fg)$ at generic level (i.e., $k$
avoiding the critical level and the admissible rational values),
the chiral algebra is Koszul (Volume~I, Theorem~A) and the PBW
spectral sequence degenerates (Volume~I, MC1).  Volume~I's
Theorem~H then gives the polynomial ring structure
$\ChirHoch^*(\mathcal{W}^k(\fg)) \cong
\mathbb{C}[\Theta_1, \ldots, \Theta_r]$, where $r = \operatorname{rank}(\fg)$
and the generators $\Theta_i$ have degrees $m_i + 1$
corresponding to the exponents of~$\fg$.
\end{proof}
